\setlength{\footskip}{8mm}

\chapter{Conclusion and Recommendations}
\label{ch:conclusion}

\section{Conclusion}
\label{sec:conclusion-conclusion}

This thesis addressed a specific gap identified in
Chapter~\ref{ch:literature-review}: existing Transformer-based methods
for monocular 3D human pose estimation scale quadratically with sequence
length, and existing state-space-based alternatives, while linear in
sequence length, order the skeleton's joints either by an anatomically
arbitrary fixed index or by a learned scan grouping, rather than by the
skeleton's own kinematic structure (Section~\ref{sec:lit-summary-gap}).

Chapter~\ref{ch:methodology} presented KinecMamba, a dual-stack, Mamba-based
architecture for this task, built around a four-directional cross-scan
selective state space layer (Section~\ref{sec:method-bistssm}), and
introduced this thesis's central contribution: a breadth-first search
kinematic-tree scan order (Section~\ref{sec:method-bfs}) that replaces
the naive joint-index order with a traversal derived directly from the
skeleton's parent-child hierarchy, together with an anatomy-aware
bone-length loss that reinforces the same skeletal structure.
Chapter~\ref{ch:results} set out the evaluation of KinecMamba on Human3.6M,
which places it in the same efficiency tier as the smallest baselines
compared against, namely PoseMamba-S and SasMamba. At roughly 1.26M
parameters it reaches 41.55mm MPJPE and 34.52mm P-MPJPE. Measured against
PoseMamba-S, the closest baseline in both architecture and parameter budget
(41.8mm / 35.0mm), this is an improvement on both metrics without an increase
in model capacity, and it is competitive with the other compact state-space
model, SasMamba. The cross-dataset evaluation on MPI-INF-3DHP and the
scan-order ablation were left as future work.

The Human3.6M result supports the thesis's motivating hypothesis, that a
scan order derived from the skeleton's kinematic tree can improve a
state-space-based pose estimator without increasing model capacity: at a
comparable parameter budget KinecMamba lowers both MPJPE and P-MPJPE relative
to PoseMamba-S. This support should be read with one caveat. As discussed in
Section~\ref{sec:results-ablation}, a comparison on Human3.6M alone does not
by itself isolate the BFS scan order as the specific cause of the difference
from PoseMamba-S, as opposed to incidental differences in training
configuration. The scan-order ablation designed to make that attribution
directly, and the MPI-INF-3DHP cross-dataset evaluation committed to in
Section~\ref{sec:intro-objectives}, were left as future work
(Sections~\ref{sec:results-mpiinf3dhp}~and~\ref{sec:results-ablation}). A
claim that the BFS kinematic-tree scan order specifically, rather than the
architecture or training configuration as a whole, drives the improvement
should therefore be treated as provisional pending that ablation.

\section{Recommendations}
\label{sec:conclusion-recommendations}

The following work is recommended to complete and extend this thesis:

\begin{enumerate}
  \item \textbf{Complete the pending evaluations.} The scan-order
    ablation study (Section~\ref{sec:results-ablation}) and the
    MPI-INF-3DHP evaluation (Section~\ref{sec:results-mpiinf3dhp}) should
    be finished and reported, and the training hyperparameters and
    hardware configuration omitted from Section~\ref{sec:results-implementation}
    should be finalized and documented, so that the results in this
    thesis are fully reproducible.

  \item \textbf{Extend the ablation beyond a single model size and sequence
    length.} Because the Human3.6M comparison in Section~\ref{sec:results-h36m}
    is conducted at a single, small parameter budget and a fixed 243-frame
    window, it remains untested whether the benefit of the BFS kinematic-tree
    scan persists, grows, or diminishes at larger model capacities such as
    those of PoseMamba-L or MotionAGFormer-L, or at longer input windows,
    where a state space scan's implicit memory can begin to decay
    \cite{zhang2024kmm}.

  \item \textbf{Explore hybrid fixed/learned scan orders.} The BFS scan
    order is fixed by the skeleton's kinematic tree
    (Section~\ref{sec:method-bfs}), while methods such as
    \cite{nguyen2025skeletonmamba} learn scan groupings from data. A
    natural extension is a hybrid approach that initializes a learned scan
    from the BFS kinematic-tree order rather than from a random or
    index-based starting point, potentially combining the interpretability
    of a kinematically-grounded order with the flexibility of a learned
    one.

  \item \textbf{Test generalization beyond the evaluated scope.} As noted
    in Section~\ref{sec:intro-limitations}, this thesis is restricted to
    single-person, monocular, fixed-topology input. Extending the BFS
    kinematic-tree scan to multi-person scenes, to skeletal topologies
    other than the Human3.6M 17-joint format (e.g.\ hand or face
    keypoints), or to in-the-wild video without ground-truth 3D annotation
    would test how specific the method's benefit is to the evaluated
    setting.

  \item \textbf{Report qualitative and per-action results.} Aggregate
    MPJPE and P-MPJPE (Section~\ref{sec:results-h36m}) do not indicate
    which actions or motion types benefit most from the proposed scan
    order. A per-action breakdown and qualitative visualization of
    predicted versus ground-truth motion sequences would clarify where the
    architectural contribution of this thesis helps most.
\end{enumerate}

\FloatBarrier

